\documentclass[12pt]{article}
\usepackage[utf8]{inputenc}
\usepackage[T1]{fontenc}
\usepackage{hyperref}
\title{Vaagri Booli Language Resources for K-12 Curriculum Development}
\author{Compiled by Claude Sonnet 5 (claude-sonnet-5) for Andr\'as Kornai}
\date{}
\begin{document}
\maketitle
\section{Introduction}
Vaagri Booli (ISO 639-3: vaa) is an endangered Indo-Aryan language of southern India, spoken by the Narikuravar/Vaagri community and showing Tamil and Kannada influence. Web search found very little material; this report lists only what appeared in search results and should be treated as a starting point.
\section{Primary Educational Resources}
No primers, textbooks or school programmes were found in search results.
\section{Linguistic and Reference Materials}
\begin{itemize}
\item Vaagri Booli language (Wikipedia): overview, phonology, and a citation of G. S. Varma, \emph{Vaagri Boli: An Indo-Aryan Language} (Annamalai University, 1970), the only full-length study named; its availability online was not confirmed. \url{https://en.wikipedia.org/wiki/Vaagri_Booli_language}
\item Bharatavani language portal lists a multilingual Vaagri boli dictionary (Vaagri Booli, Tamil, Hindi, Gujarati, English); no publisher or date was given in the listing. \url{https://bharatavani.in/gujarati/bhashakosha?page=3}
\item Numerals data for Vaagri Booli (Max Planck Institute Lingweb), small dataset, not a grammar. \url{https://lingweb.eva.mpg.de/channumerals/Vaagri-Booli.htm}
\end{itemize}
Caution: internationalphoneticalphabet.org pages on Vaagri Booli and Vaghri appeared low quality and partly conflict with scholarly sources; they are not recommended. The separate Vaghri (Gujarat) language should not be confused with Vaagri Booli.
\section{Religious and Community Resources}
None found in search results.
\section{Digital and Audio Resources}
None found in search results.
\section{Curriculum Development Context}
The language appears to be unwritten in school contexts and not taught in schools. Curriculum work would need fieldwork with the community and linguists, starting from the Varma study and the Bharatavani dictionary.
\end{document}